% ======================= 第 一 章 =======================
\chapter{月亮上的古华夏}
\chaptersubtitle{ГЛАВА ПЕРВАЯ}{ЛУННЫЙ КИТАЙ: ПРАРОДИНА В МОРЕ СПОКОЙСТВИЯ}
\epigraphbox{“遂古之初，谁传道之？上下未形，何由考之？”}{——屈原《天问》}

\noindent
迄今为止，一切关于中国文明起源的讨论，都建立在同一个未经检验的前提之上：文明的演进是连续的。人们默认，黄河流域的村落、城邑、文字与礼制，是从更早的村落、城邑、文字与礼制中一步步生长出来的；其间纵有战乱与饥馑，脉络从未真正中断。

这一前提从未被证明。它之所以被普遍接受，仅仅因为——如本书将要指出的——\kw{中断的证据被系统地清除了，而清除者正是中断之后重建文明的那批人。}

本章所要论证的命题可以概括为一句话：中华文明并非从黄河流域的土地上生长出来，而是从月球静海的一片穹顶之下，被一艘名为“昆仑山号”的亚光速方舟，连同其全部技术、制度与记忆，一并投放在黄河流域的。此后三千年间，这批移民因大气适应性退化而逐代丧失其技术能力，最终把自己的来历改写为神话，把自己的仪器改写为礼器，把自己的失败改写为天命。

这个命题在方法上是可检验的。检验它的工具不是训诂，而是质谱。

\section{生态重启：一种被忽略的周期律}

一九五七年，苏联科学院第四考察队在贝加尔湖南端取得一根长十一米的湖底沉积岩芯。岩芯在 4.71 米与 6.03 米处各出现一层厚约两厘米的白色火山灰夹层，其铱元素丰度分别为背景值的 41 倍与 37 倍。按当时的年代学，这两层的年龄分别被定为距今一万一千二百年与两万四千六百年。\footnote{Б. П. Соколов, \textit{Осадочные керны Байкала}, М., 1961, с. 88—92. 该报告第七章至第九章在一九六四年的英译本中被整章删除，删除理由在译本前言中写作“技术性调整”。}

真正令巴兰尼科夫教授不安的，不是铱异常本身——那只是撞击事件的常规证据——而是夹层之上三十厘米内，花粉谱中\kw{禾本科与蒿属的比例发生了不可逆的翻转}，而这一翻转在所有已知的第四纪记录中都需要至少六千年才能完成。换言之，沉积序列在这一层之上“跳”过了六千年。

这一现象，他在一九七三年的一篇短文里首次命名为“生态重启”（\ru{экологи\-ческий перезапуск}）：\kw{当一个文明的资源消耗超过其所在生态圈的恢复阈值时，维持该生态圈的技术系统会主动执行一次全局回滚，把生物圈恢复到某个较早的稳定态，同时把文明的技术存量削减到该稳定态所能承载的水平。}\footnote{А. В. Баранников, «Экологический перезапуск как историческая категория», \textit{Вопросы истории}, 1973, № 9, с. 44—61. 该文发表后，作者被要求“在方法上更加谨慎”，此后二十年未再使用这一术语，直至一九九三年才在本书俄文第一版中恢复。}

必须指出，这一假说在方法论上的优势是决定性的：它不需要假设任何超自然力量，只需要假设一个极其古老、极其耐心、且至今仍在运行的\kw{行星级生态维持系统}。至于该系统由谁建造、位于何处，属于本书第三章的讨论范围；本章只需证明该系统曾经运行过，并且至少运行过四次。

\begin{table}[htbp]
\centering
\caption{四次生态重启的分期与特征（据巴兰尼科夫体系，2019 年修订）}
\label{tab:reboot}
\small
\begin{tabular}{@{}clllr@{}}
\toprule
\textbf{序次} & \textbf{时间（距今）} & \textbf{触发因素} & \textbf{技术存量削减} & \textbf{持续时间}\\
\midrule
R-1 & 约 21{,}300 年 & 地核热流异常 & 97.4\% & 1{,}100 年\\
R-2 & 约 11{,}600 年 & 冰盖—海洋耦合失稳 & 99.1\% & 2{,}400 年\\
R-3 & 约 5{,}100 年 & 大气氧分压漂移 & 94.8\% & 900 年\\
R-4 & 约 2{,}200 年 & 轨道工业残骸再入 & 99.982\% & 3{,}300 年\\
\bottomrule
\end{tabular}
\tabsource{巴兰尼科夫《重构史学导论》，莫斯科，2015 年，第 118 页表 4。R-4 的“持续时间”一栏有作者手书补注：“尚未结束。”
}
\end{table}

表 \ref{tab:reboot} 中最值得注意的不是削减幅度，而是 R-4 的持续时间。按巴兰尼科夫的算法，R-4 开始于公元前二百二十一年，其回滚过程至今仍未完成——我们正生活在一次生态重启的中段。这一判断的直观后果是：\kw{我们今天所见的一切“古代”，都是重启之后的重建物，而非重启之前的原物。}考古学所揭示的，从来不是文明的起源，而是文明废墟上长出的第二茬作物。

\begin{archive}[绝密／СОВЕРШЕННО СЕКРЕТНО]{档案 № ФВ-1963-07}
本卷宗系一九六三年九月由苏联科学院远东委员会第四组移交，共 217 叶。内容为一九五九年至一九六二年在蒙古人民共和国肯特省、中国内蒙古自治区鄂尔多斯高原及苏联赤塔州三地采集的“异常金属样品”分析记录。

第 88 叶记载：鄂尔多斯样品 Э-14 为一枚重 41.6 克的片状物，成分分析显示为 \textsc{nb}—\textsc{ta}—\textsc{w} 三元合金。其铌同位素比值 \textsc{nb}-94/\textsc{nb}-93 为 0.0172，与地球物质标准值 0.0170 在误差范围内一致；但 \textsc{ta}-180 的丰度较地球标准值高出 2.3 倍。

第 89 叶页边有铅笔批注，字迹经鉴定为巴兰尼科夫本人：“坦塔洛斯之泪。他们把它带下来了。”该批注下方另有一行被涂黑，经红外成像可辨出“重量不对”四字。
\end{archive}

\section{广寒宫：编号 \textsc{mare-tranq-07}}

“广寒宫”并非文学想象，而是一个工程编号的中文雅化。在莫斯科档案馆第 117 全宗的原始图纸上，该设施的单位代号为 \textsc{mare-tranq-07}，即“静海七号”。图纸现存三张，比例尺分别为 1:2000、1:200 与 1:20；第三张仅存右下角，绘有一座环形结构的局部剖面，标注“第 4 约束段，检修周期 4000 h”。

复原工作依据三条彼此独立的线索。

其一，月面正面的低反射率异常区。静海西缘存在一片直径约 3.2 公里的异常平滑区，其雷达反照率标准差仅为周围月海的七分之一。行星地质学将此解释为“年轻熔岩流”，但熔岩流不会呈现\kw{正八边形}的边界。\footnote{该异常区编号 \textsc{mt}-07，见美国地质调查局一九七九年出版的 1:250000 月面地质图（I-1234）第 3 幅。巴兰尼科夫在此图复印件上以红笔描出八边形边界，旁注：“玄武岩不会遵守建筑规范。”}

其二，轨道残骸层。地球同步轨道上有一类“来源不明”的微小碎片，其成分以钛—铝—钒合金为主，表面有周期性微陨石撞击坑，坑密度对应的暴露年龄约为两万年。这类碎片在任何国家的发射记录中都不存在。\footnote{参见欧洲空间局《空间环境报告》2019 年版，第 6 章“未编目物体”。该章作者在结论中写道：“它们的轨道倾角分布过于整齐，整齐得不像事故的产物。”}

其三，也是最关键的一条，是\kw{月壤中的氦-3 富集层}。静海西部一处直径八百米的区域内，氦-3 丰度较周围高出一个数量级，且富集层边界呈规则的同心环状。这只能是长期、定向、可控的排气沉积所致。\footnote{А. В. Баранников, \textit{Лунная прародина: опыт реконструкции}, М., 2018, с. 203—217. 作者在此处计算了所需排气量，结论是“约等于一座百万人口城市四百年的呼吸总量”，并补充：“他们呼吸得很均匀。”}

\begin{figure}[htbp]
\centering
\begin{tikzpicture}[scale=1.02, font=\footnotesize]
  \useasboundingbox (-4.3,-4.3) rectangle (7.9,4.3);
  % 外墙与栅格
  \draw[line width=0.9pt] (0,0) circle (3.5);
  \draw[line width=0.4pt, dash pattern=on 3pt off 2pt] (0,0) circle (3.05);
  % 同心环
  \foreach \r in {2.45,1.85,1.25}{
    \draw[line width=0.5pt] (0,0) circle (\r);}
  % 正八边形边界
  \draw[line width=0.5pt, dash pattern=on 1pt off 1.4pt]
    (0:3.5) -- (45:3.5) -- (90:3.5) -- (135:3.5) -- (180:3.5)
    -- (225:3.5) -- (270:3.5) -- (315:3.5) -- cycle;
  % 中央枢纽
  \fill[black!10] (0,0) circle (0.42);
  \draw[line width=0.7pt] (0,0) circle (0.42);
  \draw[line width=0.4pt] (0,0) circle (0.2);
  \fill (0,0) circle (0.05);
  % 六个居住舱
  \foreach \a in {30,90,150,210,270,330}{
    \begin{scope}[rotate=\a]
      \draw[line width=0.6pt, fill=black!5] (-0.32,2.12) rectangle (0.32,2.72);
    \end{scope}}
  % 着陆坪
  \fill[pattern=north east lines, pattern color=black!45]
    (1.05,-1.05) -- (2.15,-0.35) -- (2.95,-1.55) -- (1.85,-2.25) -- cycle;
  \draw[line width=0.6pt]
    (1.05,-1.05) -- (2.15,-0.35) -- (2.95,-1.55) -- (1.85,-2.25) -- cycle;
  % 引线
  \draw[line width=0.3pt] (0.42,0.30) -- (3.05,2.55) -- (3.75,2.55);
  \node[anchor=west] at (3.78,2.55) {中央枢纽（\textsc{core}）};
  \draw[line width=0.3pt] (2.42,2.35) -- (3.05,1.40) -- (3.75,1.40);
  \node[anchor=west] at (3.78,1.40) {居住环 \textsc{i}—\textsc{vi}};
  \draw[line width=0.3pt] (2.62,-1.30) -- (3.75,-0.75);
  \node[anchor=west] at (3.78,-0.75) {着陆坪 \textsc{lp-3}};
  \draw[line width=0.3pt] (-3.40,-1.00) -- (-3.95,-1.00);
  \node[anchor=east] at (-3.98,-1.00) {外层屏蔽环};
  % 指北针与比例尺
  \draw[line width=0.4pt, -{Latex[length=1.6mm,width=1.2mm]}] (-3.95,2.55) -- (-3.95,3.30);
  \node[font=\scriptsize, anchor=south] at (-3.95,3.32) {N};
  \draw[line width=0.5pt] (-4.05,-3.60) -- (-3.05,-3.60);
  \draw[line width=0.4pt] (-4.05,-3.68) -- (-4.05,-3.52);
  \draw[line width=0.4pt] (-3.05,-3.68) -- (-3.05,-3.52);
  \node[font=\scriptsize, anchor=north] at (-3.55,-3.74) {500 м};
\end{tikzpicture}
\caption{静海七号（“广寒宫”）基地总平面复原示意图}
\label{fig:guanghan}
\figsource{据莫斯科档案馆第 117 全宗第 4 目第 2187 号卷宗第 12 叶重绘，比例尺 1:2000。另见巴兰尼科夫《月面设施复原图集》，圣彼得堡，2016 年，图版 11。原图右下角盖有“已校对”印章，日期为 1961 年 4 月 7 日。
}
\end{figure}

图 \ref{fig:guanghan} 所复原的，是一个可容纳十一万四千人的封闭生态体。其设计逻辑之严密，足以令任何一位现代空间站工程师感到羞愧：六个居住环按人口密度递减排布，中央枢纽兼作热控与姿态调整，着陆坪置于下风向，以避免推进剂残留污染居住区。

而最令笔者（以及译者）震动的，是图纸左上角的一行小字。俄文，铅笔，写法极其潦草：

\vspace{2mm}
\begin{center}
{\itshape\small «Материковый вариант — не проверен».}
\end{center}
\begin{center}
{\small\color{pencil}（“大陆方案——未验证”。）}
\end{center}
\vspace{2mm}

这行批注写于公元一九六一年。写下它的人，显然已经知道后来发生了什么。

\section{昆仑山号}

公元前二二一年，秦将王贲引黄河水灌大梁，魏亡。同年，秦始皇并天下，书同文，车同轨。

传统史学把这一年视为中国历史上第一个统一帝国的起点。本书不否认这一点，但必须补充：\kw{在同一个年份，一个远比帝国更重要的事件正在大气层外发生——广寒宫的最后一批居民，启动了撤离程序。}

撤离的原因，巴兰尼科夫给出的解释冷静得近乎残忍：\kw{地核冷却。}月面基地的能源、热控与生保系统，本质上是地球行星发动机的附属设施；当地核的热对流效率在约二百年间下降 3.1\% 时，月面设施所依赖的定向能输送链路便无法维持。这是一个纯工程问题，与道德无关。

于是有了“昆仑山号”。

关于这艘方舟，我们所知的一切都来自第二手材料，因为船体本身仍在地球轨道上——以“月球”的名义。\footnote{此说初看荒谬，但请注意：月球自转周期与公转周期严格相等（潮汐锁定），其轨道偏心率仅 0.0549，且月壳平均厚度（约 35 公里）远大于类地行星的正常值。巴兰尼科夫在《月面设施复原图集》第 7 页写道：“一个天体若同时满足这三个条件，那它要么是奇迹，要么是船。”}

\begin{figure}[htbp]
\centering
\begin{tikzpicture}[scale=0.95, font=\small]
  \begin{axis}[
    width=13.0cm, height=6.0cm,
    xlabel={飞行时间（小时）}, ylabel={},
    xmin=0, xmax=264, ymin=0, ymax=12.5,
    axis lines=left, tick align=outside,
    xtick={0,48,96,144,192,240},
    ytick={0,2,4,6,8,10,12},
    every axis label/.style={font=\footnotesize},
    tick label style={font=\scriptsize},
    grid=major, grid style={rulegrey!45, line width=0.2pt},
    ]
    \addplot[line width=0.9pt, domain=0:264, samples=120]
      {12.1*exp(-x/90)+0.55};
    \addplot[line width=0.5pt, dashed, domain=0:264, samples=120]
      {9.4*exp(-x/74)+0.42};
    \node[font=\scriptsize, anchor=west] at (axis cs:150,6.4) {方案 A（静海直降）};
    \node[font=\scriptsize, anchor=west] at (axis cs:150,3.1) {方案 B（近地制动）};
    \draw[line width=0.3pt, dashed] (axis cs:0,0.55) -- (axis cs:264,0.55);
    \node[font=\scriptsize, anchor=north west] at (axis cs:14,0.42) {末速 0.55 km/s};
    \node[rotate=90, font=\footnotesize] at (rel axis cs:-0.115,0.5) {速度（km/s）};
  \end{axis}
\end{tikzpicture}
\caption{“昆仑山号”再入速度衰减曲线（两条候选方案）}
\label{fig:traj}
\figsource{据档案第 117 全宗第 2187 号卷宗第 41 叶数据点重绘。原图纵轴无标注，数值由俄文标记“в. е.”换算而来，换算系数见巴兰尼科夫 2018 年书第 244 页。}
\end{figure}

图 \ref{fig:traj} 中方案 A 与方案 B 之差，是本书全部结论的力学基础。方案 A 保持较高速度穿越大气，需要 4.7 公里长的船体作为减速面；方案 B 则在近地轨道完成制动，以较低速度进入，但要求船体具备完整的再入结构。档案在两条曲线旁各有一处批注：方案 A 旁写“已执行”，方案 B 旁写“未执行”。

既已执行，则我们必须接受它的全部后果：一艘四公里长的物体，以高超音速，在公元前三世纪的一个早晨，落入黄河流域。

\begin{table}[htbp]
\centering
\caption{“昆仑山号”降落参数复原（方案 A）}
\label{tab:landing}
\footnotesize
\begin{tabular}{@{}llr@{}}
\toprule
\textbf{参数} & \textbf{数值} & \textbf{不确定度}\\
\midrule
船体长度 & 4{,}713 м & ± 12 м\\
再入质量 & 8.4 × 10\textsuperscript{7} t & —\\
大气进入角 & 6.2° & ± 0.3°\\
峰值热流 & 1.47 × 10\textsuperscript{4} kW/m² & ± 4\%\\
落点 & 34.6°N, 110.4°E & ± 40 км\\
撞击当量 & 2.1 × 10\textsuperscript{3} Mt & —\\
抛射物总量 & 1.9 × 10\textsuperscript{12} t & —\\
\bottomrule
\end{tabular}
\tabsource{巴兰尼科夫《月面设施复原图集》，第 251 页。落点坐标 34.6°N、110.4°E 位于今河南省灵宝市以西，恰在仰韶文化西缘与二里头文化之间的过渡带上。作者在此处加注：“他们挑了一个好地方。”
}
\end{table}

表 \ref{tab:landing} 中“抛射物总量”一栏，是理解中国北方第四纪晚期地层的关键。一次 2.1×10\textsuperscript{3} 百万吨级的撞击，会把约一万九千亿吨物质抛入平流层。这些物质在随后的三百年里缓慢沉降，形成我们今日在黄土高原上所见的、厚达数十米的\kw{次生黄土}。地质学界一向将其解释为“风力搬运”，此说并不错误——只是搬运工不是风，是冲击波。

\section{大气适应性退化与“绝地天通”}

落点选得好，不等于活得好。

方舟带来的十二万人，是在月面低重力、封闭大气、全人工光照条件下繁衍了数千代的种群。他们的血红蛋白氧亲和力、肺泡表面积与皮肤黑色素分布，都是为月球环境优化的。将其投入地球的一点二个大气压、百分之二十一氧分压、含尘量极高的环境中，后果不言而喻。

档案对此有极简的记录。第 2187 号卷宗第 96 叶是一张表格，标题为“第二代存活率”。数据如下：

\begin{center}
\footnotesize
\begin{tabular}{@{}lcccc@{}}
\toprule
世代 & G1 & G2 & G3 & G4\\
\midrule
存活至成年（\%） & 61.4 & 38.7 & 12.9 & 4.2\\
\bottomrule
\end{tabular}
\end{center}

第四代之后，表格中断。第 97 叶只有一行字，仍是铅笔：“他们开始挖洞。”

这行字与中国考古学的全部发现严丝合缝。仰韶文化的半地穴式房屋、龙山文化的窑洞、二里头宫殿的夯土基址——所有这些“建筑成就”，在本书的框架内都是同一件事的不同阶段：\kw{一个正在窒息的种群，用尽最后的智力，为自己挖掘能够过滤空气的掩体。}

而技术，正是在这个过程中被逐代丢弃的。这不是遗忘，而是\kw{注意力的重新分配}：当一个族群的全部认知资源都必须用于维持呼吸时，它没有余力维持一台质谱仪。巴兰尼科夫将此称为“注意力坍缩”，并给出了一个冷酷的量化模型：

\begin{equation}
\frac{dT}{dt} = -\kappa \, T \left(1 - \frac{E_{\mathrm{s}}}{E_{\mathrm{t}}}\right)
\label{eq:collapse}
\end{equation}

式中 $T$ 为技术存量，$E_{\mathrm{s}}$ 为维持生存所需的能量，$E_{\mathrm{t}}$ 为可获得的总能量。当生存耗能占比趋近于 1 时，技术存量的衰减率趋近于零——不是因为技术被保住了，而是因为\kw{已经没有什么技术可以再失去了}。\footnote{式 \eqref{eq:collapse} 的形式与洛特卡—沃尔泰拉方程同构。作者本人从未提及这一点，但他在 1993 年的一封信中承认：“我只是把捕食者换成了呼吸。”}

于是我们来到了中国上古史上最著名、也最被误读的一句话。

《国语·楚语下》载观射父之言：“及少皞之衰也，九黎乱德，民神杂糅，不可方物……颛顼受之，乃命南正重司天以属神，命火正黎司地以属民，使复旧常，无相侵渎，是谓绝地天通。”

历代注疏家都将“绝地天通”解释为宗教改革，即垄断通天之权。这一解释在文本上无懈可击，在物理上却荒谬：\kw{如果“天”只是一组抽象的神祇，何须“绝”之？}

“绝地天通”的字面义是：切断地与天之间的通道。而在本书的框架内，它只有一个含义——\kw{关闭轨道电梯。}

第 2187 号卷宗第 121 叶，是一份用打字机誊抄的“设施退役清单”，共十九项。其中第七项写作：

\begin{center}
\footnotesize\itshape
“7. Вертикальный транспорт (лифт) — демонтирован.\\
Опоры срезаны на высоте 40 м. Маркировка на срезе: «НЕ КОПАТЬ».”
\end{center}
\begin{center}
\small\color{pencil}
（第 7 项：垂直运输装置（电梯）——已拆除。\\
支柱于 40 米高度截断。截断处标记：“勿挖”。）
\end{center}

“勿挖”两个字，用的是中文。

\begin{archive}[绝密／СОВЕРШЕННО СЕКРЕТНО]{档案 № ФВ-1971-22}
一九七一年，中国科学院考古研究所在河南偃师二里头遗址第三发掘区，于 T-14 探方距地表 4.2 米处发现一处“柱洞群”。共九个柱洞，呈正圆形排列，直径 3.1 米；洞内填充物为经高温烧结的土块，烧结温度经热释光测定为 1{,}240 ℃ ± 60 ℃。

发掘报告判定此处为“商代早期大型建筑基址的柱础”。\footnote{《河南偃师二里头遗址三区发掘简报》，《考古》1974 年第 5 期，第 268—279 页。报告第 273 页有一句后被删除的话：“柱洞底部发现不明金属残片，已送有关部门鉴定。”此后未见鉴定结论发表。}

同层出土陶片七件，其中一件（编号 71YLT14:6）内壁有朱书符号两个。符号经红外成像可辨为“天”“梯”二字，字迹草率，笔画末端有拖曳痕迹，显系仓促写就。
\end{archive}

我必须在这里插入一句方法论上的说明，尽管它可能会令部分读者不适。

上述材料，单看每一条，都可以有别的解释：铱异常可以是撞击，柱洞可以是柱础，朱书可以是随手涂鸦。本书的方法不是证明每一条都是铁证，而是\kw{证明这些解释的集合在统计上是不可能的}。当十七个彼此独立的异常，在同一地层、同一区域、同一时间窗口内同时出现时，坚持“每一条都有平凡解释”所需假设的总量，已经超过了接受一个不平凡解释所需的假设量。

这就是奥卡姆剃刀在我手中的用法：\kw{我不削减实体，我数实体。}

\section{本章结论}

公元前二二一年之后的三千年里，黄河流域的居民重建了一个文明。他们用青铜代替了合金，用礼器代替了仪器，用天命代替了工程手册。他们把祖先的降落称为“天命玄鸟，降而生商”，把祖先的穹顶称为“天圆地方”，把祖先的雪山称作“昆仑”，把祖先的失败称为“绝地天通”。

这些改写并非欺骗，而是\kw{一种在认知资源极度匮乏条件下的最优编码}：一个必须让后代记住“不要挖那里”的族群，最经济的办法就是把那条禁令写进神话，写进仪式，写进每一次祭祀的方位与次序。

于是“勿挖”变成了“不可亵渎”，“轨道电梯”变成了“建木”，“月面基地”变成了“广寒宫”。

神话不是历史的对立面。神话是历史在记忆带宽不足时的压缩格式。\footnote{这一表述首次出现在巴兰尼科夫 1979 年的笔记中，后因“措辞不当”被苏联时期出版社要求改写为“神话是历史的口头传承形式”。本书第三版恢复了原始措辞。}

而我们的任务，是解压缩。
